\documentclass[14pt]{beamer}

% Assets
\usepackage[czech]{babel}				% Jazyk
\usepackage[a-2u]{pdfx}					% Kopírování z pdfka
\usepackage{tikz}						% Schémata automatů
\usepackage{csquotes}					% české uvozovky
\usepackage{enumerate}					% enumerate environment
\usepackage{indentfirst}
\usepackage{mathtools}
\usepackage{pifont}
\usepackage{xcolor}
\usepackage{enumitem,xcolor}
\usepackage{amsmath}
\usepackage[utf8]{inputenc}

\usepackage{listings}                   % Úryvky z kódu (C#)
\lstset{language=[Sharp]C, frame=lr}
% Beamer theme
\usetheme{Boadilla}
\setbeamertemplate{frame numbering}[fraction]
\usecolortheme[named=black]{structure}
\setbeamertemplate{navigation symbols}{}
\setbeamerfont{title}{series=\bfseries,parent=structure}
\setbeamerfont{frametitle}{series=\bfseries,parent=structure}
\usefonttheme[onlymath]{serif}
\urlstyle{same}

% Dark theme
\setbeamercolor{frametitle}{fg=white}
\setbeamercolor{title}{fg=white}
\setbeamercolor{background canvas}{bg=black}
\setbeamercolor{normal text}{fg=white}

\defbeamertemplate*{title page}{customized}[1][]
{ 
  \usebeamerfont*{title}\inserttitle\par
  \bigskip
  \usebeamerfont*{subtitle}\textit{\insertsubtitle}\par
  \bigskip \bigskip \bigskip \bigskip 
  \usebeamerfont{author}\insertauthor\par
  \usebeamerfont{institute}Kabinet \office\\\url{weber3@spsejecna.cz}\bigskip
}
% Enumerate
%\setlist[enumerate]{topsep=0pt,itemsep=-1ex,partopsep=1ex,parsep=1ex,label=(\arabic*)}

\MakeOuterQuote{"}

% Colors
\definecolor{lightblue}{HTML}{009AD4}
\definecolor{darkgreen}{HTML}{0D7103}
\definecolor{lightgreen}{HTML}{68FF00}
\definecolor{darkred}{HTML}{AF0B0B}
\definecolor{lightred}{HTML}{FF5100}
\definecolor{orange}{HTML}{FFE000}
\definecolor{codeblue}{HTML}{FF0055}
\definecolor{codegreen}{rgb}{0,0.6,0}
\definecolor{codegray}{rgb}{0.5,0.5,0.5}
\definecolor{codebeige}{HTML}{D4A000}
\definecolor{backcolour}{rgb}{0.95,0.95,0.92}

\newcommand{\markred}[1]{\textcolor{lightred}{#1}}
\newcommand{\markgreen}[1]{\textcolor{lightgreen}{#1}}
\newcommand{\markorange}[1]{\textcolor{orange}{#1}}
\newcommand{\markblue}[1]{\textcolor{lightblue}{#1}}

% Inline images
\newcommand{\inlineimgscale}{1.1}

% X and check mark
\newcommand{\cmark}{\markgreen{\ding{51}}}
\newcommand{\xmark}{\markred{\ding{55}}}

% Redefinions
\renewcommand{\implies}{\Rightarrow}
\renewcommand{\impliedby}{\Leftarrow}

% Math
\newcommand{\R}{\mathbb{R}}
\newcommand{\C}{\mathbb{C}}
\newcommand{\N}{\mathbb{N}}
\newcommand{\Z}{\mathbb{Z}}
\newcommand{\Q}{\mathbb{Q}}

% Code
\lstdefinestyle{clang}{
    basicstyle=\small\ttfamily\color{white},
    language=C,
    keywordstyle=\color{codeblue},
    commentstyle=\color{codegreen},
    numberstyle=\tiny\color{codegray},
    stringstyle=\color{codebeige},
    breakatwhitespace=false,
    breaklines=true,
    captionpos=b,
    keepspaces=true,
    numbersep=5pt,
    showspaces=false,
    showstringspaces=false,
    showtabs=false,
    morekeywords={void,int,double,float,unsigned,if,else,\#include}
    tabsize=0.5
}
\lstset{escapeinside={(*}{*)},style=clang}

\newcommand{\hlcode}[1]{\colorbox{red}{#1}}

% Title page
\subtitle{Informační a komunikační technologie}
\title{První pohled na jazyk C}
\author{David Weber}
\def\office{K13}
\def\email{weber3@spsejecna.cz}

\begin{document}

    % Itemize
    \setlist[itemize]{label=\textcolor{white}{\textbullet}}

    % Slides
    \begin{frame}
        \titlepage
    \end{frame}

    \begin{frame}[t,fragile]{Co už známe\dots}
        \begin{itemize}
            \item Tělo programu a vkládání knihoven.
            \begin{lstlisting}
#include <stdio.h>

int main(void) {

    // Kód programu

    return 0;
}
            \end{lstlisting}
            \item \texttt{printf(\dots)} -- příkaz pro výpis textového řetězce.
        \end{itemize}
    \end{frame}

    \begin{frame}[t]{Co dnes probereme}
        \begin{itemize}
            \item Čtení vstupu od uživatele
            \item Výpis proměnné
            \item Větvení programu (podmínky)
        \end{itemize}
    \end{frame}

    \begin{frame}[t]{Čtení vstupu}
        \begin{itemize}
            \item Příkaz \texttt{scanf(\textbf{<formátovací řetězec>}, \&\textbf{<proměnná>});}
            \item Formátovací řetězec
            \begin{itemize}
                \item Specifikuje \markorange{datový typ} vstupní hodnoty
                \item Základní položky:
                \begin{itemize}
                    \item \texttt{\%d} $\leftarrow$ celé číslo (\markgreen{int})
                    \item \texttt{\%f} $\leftarrow$ číslo s desetinnou částí (\markorange{float})
                    \item \texttt{\%s} $\leftarrow$ řetězec znaků (\markorange{string})
                \end{itemize}
                \item Píšeme do uvozovek \textquotedbl{} \textquotedbl (je to řetězec).
            \end{itemize}
        \end{itemize}
    \end{frame}

    \begin{frame}[t,fragile]{Příklad}
        \begin{lstlisting}
#include <stdio.h>

int main() {
    int x;
    scanf("%d", &x);
    return 0;
}
        \end{lstlisting}
    \end{frame}

    \begin{frame}[t,fragile]{Výpis proměnné}
        \begin{itemize}
            \item Provádíme také pomocí formátovacího řetězce.
            \item Pomocí \%d\,/\,\%f\,/\,\%s specifikujeme typ proměnné, kterou vypisujeme.
            \item Už nepíšeme \&.
            \begin{lstlisting}
int x;
x = 5;
printf("%d", x);
            \end{lstlisting}
            \item Program vypíše 5.
        \end{itemize}
    \end{frame}

    \begin{frame}[t,fragile]{Více hodnot}
        \begin{itemize}
            \item V rámci jednoho příkazu lze načíst/vypsat i více hodnot.
            \item Příklad pro dvě hodnoty
            \begin{lstlisting}
int x;
int y;
scanf("%d %d", &x, &y);
printf("%d %d", x, y);
            \end{lstlisting}
            \item Lze načítat/vypisovat libovolný počet hodnot.
        \end{itemize}
    \end{frame}

    \begin{frame}[t,fragile]{Podmínky}
        \begin{itemize}
            \item Umožňují reagovat na různé situace, které mohou v programu nastat.
            \item Skládá se z \textbf{logického výrazu} určujícího podmínku a tzv. \textbf{těla}.
            \item Klíčové slovo \texttt{if}.
            \begin{lstlisting}
if (podmínka) {
    // Kód programu
}
            \end{lstlisting}
        \end{itemize}
    \end{frame}

    \begin{frame}[t,fragile]{Příklad}
        \begin{lstlisting}
int x;
scanf("%d", &x);
if (x > 5) {
    printf("x je vetsi nez 5");
}
        \end{lstlisting}
    \end{frame}

    \begin{frame}[t]{Operátory}
        \begin{itemize}
            \item Seznam operátorů:
            \begin{itemize}
                \item \texttt{==} (rovnost)
                \item \texttt{!=} (nerovnost)
                \item \texttt{<} (menší než)
                \item \texttt{>} (větší než)
                \item \texttt{<=} (menší než nebo rovno)
                \item \texttt{>=} (větší než nebo rovno)
            \end{itemize}
        \end{itemize}
    \end{frame}

    \begin{frame}[t]{Logické operátory}
        \begin{itemize}
            \item Umožňují "sloučení" více podmínek do jedné.
            \item Seznam operátorů:
            \begin{itemize}
                \item \texttt{\&\&} - logický \textbf{AND} (podmínka je splněna, pokud jsou \textbf{všechny} dílčí podmínky splněny).
                \item \texttt{||} - logický \textbf{OR} (podmínka je splněna, pokud je splněna \textbf{alespoň jedna} z dílčích podmínek).
            \end{itemize}
            \item Syntaxe:
            \begin{center}
                \texttt{if (podm$_1$ \&\& podm$_2$ \&\& \dots{ }\&\& podm$_n$)}\\
                \texttt{if (podm$_1$ || podm$_2$ || \dots{ }|| podm$_n$)}
            \end{center}
        \end{itemize}
    \end{frame}

    \begin{frame}[t,fragile]{Příklad}
        \begin{itemize}
            \item Podmínka, že zadané číslo $x$ je z intervalu $\langle 0,\,5\rangle$
            \begin{lstlisting}
#include <stdio.h>

int main(void) {
    int x;
    scanf("%d", &x);

    if (x >= 0 && x <= 5) {
        printf("Cislo x je z intervalu (0, 5).");
    }
}
            \end{lstlisting}
        \end{itemize}
    \end{frame}

    \begin{frame}{Otázky?}
        \begin{figure}
            \centering
            \includegraphics[scale=.4]{images/discussion_inverted.png}
        \end{figure}
    \end{frame}

\end{document}